The broad goals that this design attempted to address are listed in Table \ref{tab:design-goals}.

\begin{table}[H]
    \centering
    \caption{Design Goals for Grasper and Detection}
    \begin{tabular}{| c || >{\centering\arraybackslash}m{0.4\textwidth} |}
        \hline
        Goal Number & Description \\
        \hline
        1 & {The device shall detect and attach to overhead, roughly cylindrical perches.}\\
        2 & {The device shall be capable of repeatedly attaching and detaching from a perch.}\\
        3 & {The device shall support a drone of at least 500g.}\\
        4 & {The mass of the device shall be under 47.5g.}\\
        5 & {The method of perch detection shall be compatible with low-cost cameras and open-source software.}\\
        6 & {The detection method shall behave deterministically, avoiding the issues noted by
        von Frankenberg and Nokleby \cite{von_frankenberg_detection_2020}.}\\
        \hline
    \end{tabular}
    \label{tab:design-goals}
\end{table}


Goals 3 and 4 were derived from surveying specifications of common commercial drones.
Table \ref{tab:drone-weights} shows the assessed models and their weights.

\begin{table}[H]
    \centering
    \caption{Specified Weight of Various Commercial Drones}
    \begin{tabular}{| c || c |}
        \hline
        {Make \& Model} & Weight [g] \\
        \hline
        DJI Mavic Mini & 249 \\
        Parrot ANAFI USA & 500 \\
        DJI Mavic Pro & 734 \\
        DJI Mavic 3 Pro & 958 \\
        DJI Mavic 3 Enterprise & 920 \\
        Autel EVO Max 4T & 1640 \\
        FLIR SIRAS & 3100 \\
        DJI Matrice 30 & 3770 \\
        \hline
    \end{tabular}
    \label{tab:drone-weights}
\end{table}


\begin{minipage}{\linewidth}
There appear to be a few classes of drones represented:
\begin{itemize}
    \item Sub 500g lightweight drones
    \item Sub 1kg medium drones
    \item Heavy weight 1.5 kg + drones
\end{itemize}
\end{minipage}
\vspace{\baselineskip}


The device must also allow for some dynamic load, particularly force introduced by motion of the perch such as a swaying tree branch.
Assuming that the drone is perched in a tree and subject to sinusoidal motion, its maximum acceleration is estimated at $63.17 m/s^2$ 
based on data provided by Kamimura et al. (0.4 m max displacement with a frequency of 2 Hz) \cite{kamimura_energy_2024}.  
The referenced study did use only Japanese Cedar, which may not
be representative of all trees. However, similar data for other species could not be found.

Equation \ref{eqn:load-lim} gives a formula for a total test mass that the device must support, 
to match the mass of a drone and its maximum dynamic load.

\begin{minipage}{\linewidth}
    \begin{equation}
        m_t = m_d (1 + \frac{a}{g}) = 7.44 m_d
        \label{eqn:load-lim}
    \end{equation}\\
    Where:
    \begin{itemize}
        \item $m_t$: Test mass to mimic the drone and its dynamic load.
        \item $m_d$: Mass of the drone.
        \item $a$: Maximum acceleration experienced by a perched drone, estimated at $63.17 m/s^2$.
        \item $g$: Gravitational acceleration, $9.81 m/s^2$.
    \end{itemize}
\end{minipage}
\vspace{\baselineskip}

This device targets the lowest weight group of drones, with mass up to 500g. Using equation \ref{eqn:load-lim},
the device must support a test load of at least 3.72 kg in order to meet goal 3.

The first reason for targeting the lowest weight group is that the low weight is far more
manageable for the device to hold. The other reasons are cost and risk. The medium and heavy drones are typically used to carry
heavy, higher quality cameras, and have a higher cost to match their larger size and higher thrust capacity
(for example the DJI Mavic Mini started at 400 USD, while the Mavic Pro started at 750 USD
\footnote{Prices as listed during the time these models were in production and sold in the US.}).
Leaving a drone perched in a tree or other location leaves it exposed to risk of damage.
Though what fraction of a researcher's budget a drone represents will vary, as will an 
individual researcher's risk tolerance, using a larger drone for extended perching may be 
particularly unappealing to the drone's owner.

Table \ref{tab:payload} shows rated payload for the models listed in Table \ref{tab:drone-weights}, which 
were used to derive goal 4.
The rated payload was calculated as the rated maximum total takeoff weight minus the drones specified weight. 
The DJI Mavic products as well as the FLIR SIRAS had no maximum takeoff weight aside from the products weight, 
so no maximum payload is listed for them. Most without a rated payload are assumed to still have some
residual lift capacity that can be used for a payload; for example the FLIR SIRAS drone has an optional
thermal camera attachment, meaning it must have additional payload capacity.

\begin{table}[H]
    \centering
    \caption{Rated Payload of Various Commercial Drones}
    \begin{tabular}{| c || c |}
        \hline
        {Make \& Model} & Payload [g] \\
        \hline
        DJI Mavic Mini & None Rated \\
        Parrot ANAFI USA & 144 \\
        DJI Mavic Pro & None Rated \\
        DJI Mavic 3 Pro & None Rated \\
        DJI Mavic 3 Enterprise & 130 \\
        Autel EVO Max 4T & 379 \\
        FLIR SIRAS & None Rated \\
        DJI Matrice 30 & 299 \\
        \hline
    \end{tabular}
    \label{tab:payload}
\end{table}

In the sub 500g category, only the Parrot ANAFI has a rated payload. This was assumed as a reasonable estimate for the
payload capacity of similarly sized drones. In order to leave some margin for safety, and possible additional task-specific
modifications, a target of no more than 33\% of that capacity, or 47.5g, was used. This target weight would also 
provide improvements over similar existing designs \cite{chi_optimized_2014,nguyen_passively_2019}.

Goal 5 is intended to keep the design accessible to a wide audience of researchers. The image processing methods 
must be selected to accommodate lower quality point cloud data, which appears to interfere with the method 
proposed by Kominami et al. \cite{kominami_detection_2023} (preliminary testing of this approach is described in 
appendix \ref{app:detect-proto}). Goal 6 aims to improve reliability compared to the method used by von Frankenberg and 
Nokleby, where the perch is sometimes not detected simply because the RANSAC method 
did not happen to randomly sample enough points from the perch \cite{von_frankenberg_detection_2020}, \cite{von_frankenberg_detection_2020}